% ==============================================================================
% PLANTILLA DE CV DE 1 PÁGINA PARA PERFILES STEM EN TRANSICIÓN A LA INDUSTRIA
% Diseñada para graduados, posdocs y científicos que buscan su primer o segundo
% rol corporativo (Data Science, Quant, Machine Learning, Decision Science).
% Compatible con Overleaf (compilar con pdfLaTeX).
% ==============================================================================

\documentclass[11pt]{article}

\usepackage{calc}
\usepackage{url}
\usepackage{amsmath}
\usepackage{lmodern}
\usepackage[T1]{fontenc}

% Márgenes optimizados para aprovechar la página sin saturarla
\usepackage[paper=letterpaper,
            top=1.3cm,
            bottom=1.3cm,
            left=1.5cm,
            right=1.5cm]{geometry}

\setlength{\parindent}{0in}
\usepackage{enumitem}
\usepackage{fancyhdr}
\pagestyle{empty} % Sin números de página (estándar para CV de 1 cuartilla)

% Enlaces e iconos limpios
\usepackage{color,hyperref}
\definecolor{darkblue}{rgb}{0.0,0.0,0.4}
\hypersetup{colorlinks,breaklinks,
            linkcolor=darkblue,urlcolor=darkblue,
            anchorcolor=darkblue,citecolor=darkblue}

\usepackage{fontawesome}

% Encabezados de sección limpios y compatibles con ATS
\newcommand{\cvsection}[1]{%
  \vspace{0.3cm}%
  {\bfseries\large\scshape #1}%
  \vspace{0.05cm}%
  \hrule height 0.6pt%
  \vspace{0.18cm}%
}

\newcommand{\semilarge}{\fontsize{14}{16.5}\selectfont}

\begin{document}

% ==============================================================================
% 1. ENCABEZADO (Directo, sin foto y SIN fecha de creación de documento)
% ==============================================================================
\begin{center}
  {\semilarge \textbf{[Tu Nombre y Apellidos]}} \\[0.1cm]
  {\large \textit{[Rol Objetivo: ej. Quantitative Researcher \textbar \ Decision Scientist \textbar \ Applied Data Scientist]}} \\[0.15cm]
  \small
  \faPhone~[+52] [Tu Teléfono] \quad \textbar \quad
  \faEnvelope~\href{mailto:tu-correo@email.com}{tu-correo@email.com} \quad \textbar \quad
  \faLinkedin~\href{https://www.linkedin.com/in/tu-perfil/}{linkedin.com/in/tu-perfil} \quad \textbar \quad
  \faGithub~\href{https://github.com/tu-usuario}{github.com/tu-usuario}
\end{center}

\vspace{-0.1cm}

% ==============================================================================
% 2. RESUMEN PROFESIONAL (3-4 líneas: quién eres analíticamente y qué resuelves)
% ==============================================================================
\cvsection{Professional Summary}
\textit{\small
[Científico cuantitativo / Doctor en Física / Ingeniero de Datos] con sólida formación en [modelado estocástico, optimización y análisis de datos a gran escala]. Especializado en traducir problemas complejos no estructurados en soluciones analíticas modulares y de alto impacto. Combina rigor científico y escepticismo metodológico con dominio de herramientas modernas de programación para acelerar la toma de decisiones basada en datos.
}

% ==============================================================================
% 3. HABILIDADES CLAVE (Visible de inmediato para superar el filtro inicial)
% [CONSEJO]: Para quienes van saliendo de academia, colocar Skills aquí demuestra
% inmediatamente que dominas herramientas productivas y no eres un perfil "solo teórico".
% Si ya tienes años en la industria, puedes mover esta sección después de Experiencia.
% ==============================================================================
\cvsection{Core Skills}
\small
\textbf{Technical \& Quantitative:} Modelado Estadístico, Optimización Matemática, Series de Tiempo, Simulación Estocástica y Monte Carlo; Machine Learning Supervisado / No Supervisado; Python (pandas, numpy, scikit-learn, scipy), SQL, Git, Linux/Bash; [C++ / Cloud: AWS / GCP si aplica].

\vspace{0.08cm}
\textbf{Business \& Strategy:} Ciencia de Decisiones, Formulación de Problemas Complejos, Traducción Negocio-Técnica, Comunicación a Stakeholders, Diseño Experimental, Metodologías Ágiles.

\vspace{0.08cm}
\textbf{Languages:} Español (Nativo); Inglés (Fluido / Profesional).

% ==============================================================================
% 4. EXPERIENCIA PROFESIONAL / INVESTIGACIÓN APLICADA
% [FÓRMULA CLAVE]: [Verbo de acción] + [Problema/Contexto] + [Impacto/Métrica]
% NOTA: El impacto se incluye dentro de cada viñeta, NO en una sección desconectada.
% ==============================================================================
\cvsection{Experience \& Applied Research}

% ROL 1: Tu rol más reciente (en industria, pasantía o estancia aplicada)
\textbf{[Cargo Objetivo / Título]} \textbar \ \textbf{[Empresa o Institución]}, \textit{[Ciudad, País]} \hfill \textit{[Año -- Presente]}
\vspace{-0.15cm}
\begin{itemize}[leftmargin=0.4cm] \small
  \item Diseñé e implementé la estrategia de análisis cuantitativo para [área o producto], identificando y priorizando oportunidades de optimización en [pricing, riesgo, segmentación o demanda].
  \vspace{-0.22cm}
  \item Traduje dinámicas complejas en modelos de [optimización / ML] modulares, acelerando el tiempo de entrega (*time-to-market*) en un [X\%] frente a la línea base previa.
  \vspace{-0.22cm}
  \item Construí pipelines de datos de punta a punta colaborando con equipos de ingeniería y negocio, asegurando que el modelo pasara a producción de forma reproducible.
\end{itemize}

\vspace{-0.15cm}

% ROL 2: Tu trayectoria de investigación científica traducida al lenguaje de la industria
% [CONSEJO]: Agrupa tu tesis doctoral o proyectos de investigación en un bloque sólido.
\textbf{[Investigador Posdoctoral / Científico Investigador]} \textbar \ \textbf{[Centro de Investigación / Univ.]} \hfill \textit{[Año -- Año]}
\vspace{-0.15cm}
\begin{itemize}[leftmargin=0.4cm] \small
  \item Dirigí el análisis cuantitativo sobre conjuntos de datos de alta dimensionalidad ($\sim$[TB / PB]) en colaboraciones científicas internacionales ([CERN / Fermilab / etc.]), aplicando técnicas de [Machine Learning / simulación Monte Carlo].
  \vspace{-0.22cm}
  \item Desarrollé algoritmos de [optimización numérica / clasificación señal-ruido] en [Python / C++], optimizando tiempos de cómputo en clústeres HPC en un [X\%] mediante paralelización.
  \vspace{-0.22cm}
  \item Coordiné y asesoré a [N] investigadores/estudiantes en modelado estadístico riguroso, control de versiones con Git y entrega de resultados bajo estándares editoriales internacionales.
\end{itemize}

% ==============================================================================
% 5. FORMACIÓN ACADÉMICA Y CREDENCIALES
% ==============================================================================
\cvsection{Education \& Credentials}
\small
\textbf{Ph.D. en [Física / Matemáticas / Computación / Ingeniería]} \textbar \ [Institución, Ciudad] \hfill \textit{[Año]} \\
\textbf{M.Sc. en [Área Científica]} \textbar \ [Institución, Ciudad] \hfill \textit{[Año]} \\
\textbf{B.Sc. en [Licenciatura / Ingeniería]} \textbar \ [Institución, Ciudad] \hfill \textit{[Año]}

\vspace{0.08cm}
\textbf{Credenciales y Distinciones:} [Certificación relevante: ej. BMV / Cloud Architect / Nivel en SNII], demostrando excelencia en resolución de problemas complejos y rigor metodológico continuo.

\end{document}
